% Test de stress de mod_ep : montée en charge jusqu'à saturation, endurance, traitements lourds
% simultanés, profil CPU et optimisations.
% Compilation : latexmk -pdf test-de-stress.tex
\documentclass[11pt,a4paper]{article}
\usepackage[french]{babel}
\usepackage[T1]{fontenc}
\usepackage[utf8]{inputenc}
\usepackage{lmodern,microtype,booktabs,tabularx,longtable,xcolor,hyperref,array}
\usepackage[margin=2.2cm]{geometry}
\hypersetup{colorlinks=true, linkcolor=blue!50!black, urlcolor=blue!50!black}
\newcommand{\code}[1]{\texttt{\detokenize{#1}}}
\newcolumntype{R}{>{\raggedleft\arraybackslash}X}
\title{Test de stress de \code{mod_ep}}
\author{Sébastien Lefebvre}
\date{Novembre 2026 -- versions 1.0.1 et 1.0.2 du module}

\begin{document}
\maketitle

\section{Objet et méthode}

Le test de charge (rapport \emph{Test de charge}) mesurait chaque page pour un seul utilisateur.
Celui-ci pousse le serveur jusqu'à la saturation, puis le maintient en charge pendant une
demi-heure, et lance en même temps les traitements les plus lourds du module. Il répond à trois
questions : où se situe la limite d'un serveur, que se passe-t-il au-delà, et le module y est-il
pour quelque chose ?

\paragraph{Serveur.} Une seule machine de 4 c\oe urs et 15~Go de mémoire, configurée comme un petit
serveur de production : nginx, PHP-FPM~8.3 (16 processus, OPcache), PostgreSQL~16 réglé
(\code{shared_buffers} 2~Go), Moodle~4.5. Le générateur de charge tourne sur la même machine et
y prend environ 17~\% du processeur : les débits mesurés sont donc un peu inférieurs à ceux d'un
serveur dédié.

\paragraph{Données.} Six promotions (une activité EP chacune, 200 EP au catalogue) : 10\,000
étudiants, 150 enseignants communs à toutes les promotions, une synthèse enseignant liée aux six
activités ; 538\,000 crédits et 475\,000 décisions en base (\code{charge_setup.php},
\code{synthese_setup.php}).

\paragraph{Charge.} Des utilisateurs virtuels (\code{stress.py}) -- 70~\% d'étudiants, 22~\%
d'enseignants, 8~\% de membres de la DEVE, répartis sur les six promotions -- ouvrent leurs pages
en boucle \textbf{sans temps de réflexion} : chacun produit la charge maximale possible. Un
utilisateur réel lit sa page une dizaine de secondes ; un utilisateur virtuel équivaut donc à
plusieurs dizaines de personnes connectées. Les étudiants s'inscrivent aussi à des EP (5~\% de
leurs requêtes : écritures sous verrou). Paliers de 25 à 400 utilisateurs, 60~s chacun.

\section{Montée en charge}

\begin{center}
\begin{tabularx}{\linewidth}{@{}r R R R R R R R@{}}
\toprule
& \multicolumn{3}{c}{\textbf{Avant optimisation}} & \multicolumn{4}{c}{\textbf{Après optimisation}} \\
\cmidrule(lr){2-4}\cmidrule(l){5-8}
\textbf{Utilisateurs} & \textbf{pages/s} & \textbf{médiane} & \textbf{99\,\%} & \textbf{pages/s} &
  \textbf{médiane} & \textbf{99\,\%} & \textbf{erreurs} \\
\midrule
25  & 37,3 & 0,65 s & 1,6 s  & 47,4 & 0,52 s & 0,95 s & 0 \\
50  & 36,5 & 1,35 s & 2,5 s  & 45,6 & 1,09 s & 1,7 s  & 0 \\
100 & 37,2 & 2,73 s & 4,7 s  & 46,8 & 2,16 s & 3,0 s  & 0 \\
200 & 37,9 & 5,58 s & 7,2 s  & 49,7 & 4,23 s & 5,0 s  & 0 \\
300 & 42,1 & 7,92 s & 11,0 s & 51,2 & 6,36 s & 9,6 s  & 0 \\
400 & 44,1 & 10,5 s & 12,0 s & 53,1 & 8,45 s & 11,3 s & 0 \\
\bottomrule
\end{tabularx}
\end{center}

\begin{itemize}
  \item \textbf{Aucune erreur} à aucun palier : ni erreur PHP ou de base de données, ni requête
    perdue, ni délai dépassé (60~s), y compris à 400 utilisateurs simultanés. Les inscriptions
    concurrentes ne produisent aucun doublon.
  \item \textbf{La saturation est atteinte dès 25 utilisateurs} sans temps de réflexion : le
    processeur est à 100~\%. Au-delà, le débit reste stable (il augmente même un peu) et les
    requêtes attendent leur tour : le temps de réponse croît proportionnellement au nombre
    d'utilisateurs, sans effondrement ni erreur. C'est le comportement attendu d'un serveur
    correctement borné (nombre fixe de processus PHP, file d'attente de nginx).
  \item La mémoire reste modeste : 2,9~Go au plus pour toute la machine à 400 utilisateurs.
\end{itemize}

\section{Où va le temps}

Un profil par échantillonnage (Excimer, \code{profil.py}) montre ce que coûte une page :

\begin{center}
\begin{tabularx}{\linewidth}{@{}X r r r r r@{}}
\toprule
\textbf{Page (après optimisation)} & \textbf{ms CPU} & \textbf{module} & \textbf{base} & \textbf{rendu} &
  \textbf{c\oe ur Moodle} \\
\midrule
Tableau de bord de l'étudiant & 73 & 3~\% & 15~\% & 26~\% & 56~\% \\
Catalogue (étudiant) & 58 & 4~\% & 12~\% & 23~\% & 61~\% \\
Validation (enseignant) & 61 & 3~\% & 11~\% & 21~\% & 64~\% \\
Participants d'un EP & 61 & 7~\% & 10~\% & 25~\% & 58~\% \\
Synthèse : à traiter (6 promotions) & 72 & 4~\% & 15~\% & 21~\% & 60~\% \\
Synthèse : validation (6 promotions) & 113 & 17~\% & 19~\% & 18~\% & 46~\% \\
Pilotage (DEVE) & 60 & 7~\% & 13~\% & 21~\% & 59~\% \\
\emph{Référence : profil utilisateur de Moodle} & \emph{41} & -- & 7~\% & 33~\% & 59~\% \\
\bottomrule
\end{tabularx}
\end{center}

Une page du module coûte de 50 à 75~ms de processeur, contre 41~ms pour une page ordinaire de
Moodle. Le code du module n'en représente que quelques pour cent : l'essentiel est le socle de
Moodle (thème, navigation, gabarits, chaînes de langue, connexion à la base), payé par toute
page. La limite d'environ 50 pages par seconde sur 4 c\oe urs est donc celle de Moodle sur
ce matériel, et non celle du module.

\section{Optimisations}

\subsection{Dans le module (version 1.0.1)}

\begin{longtable}{@{}p{5.2cm}p{10.4cm}@{}}
\toprule
\textbf{Constat} & \textbf{Correction} \\
\midrule\endhead
Places prises et états par EP : 10 à 17~ms par page de catalogue, parcours de toutes les demandes
  de la promotion. & Index (promotion, EP, état) : la requête se lit sur l'index seul, 2,6~ms. \\
File de la DEVE : 19~ms pour trier les 7\,000 demandes en attente avant d'en afficher 40 ; total
  calculé avec une jointure inutile, 17~ms. & Index (promotion, état, date) : la page se lit dans
  l'ordre de l'index, 0,3~ms ; total sans jointure, 1,5~ms. \\
Synthèse, liste des demandes de l'enseignant sur six promotions : 143~ms, dont 39~ms à mettre en
  forme le type et le nom de l'étudiant de chacune des milliers de lignes. & Libellés calculés une
  fois par type ; noms des étudiants calculés seulement pour un tri par nom ou pour les lignes
  affichées : 113~ms. \\
Participants d'un EP : toute la promotion (2\,000 inscrits) chargée à chaque affichage pour
  proposer l'ajout d'étudiants. & Candidats lus en base avec la recherche et une limite : 83 à
  61~ms. \\
Libellés d'état et d'année d'étude recalculés à chaque ligne (sept chaînes par ligne pour
  l'année). & Calculés une fois par requête. \\
\bottomrule
\end{longtable}

\subsection{Dans le déploiement}

\begin{itemize}
  \item \textbf{Connexions persistantes à la base} (\code{$CFG->dboptions['dbpersist'] = 1}, ou un
    pool de connexions de type pgbouncer) : la connexion à PostgreSQL coûtait environ 9~\% du
    processeur de chaque page ; à 100 utilisateurs, le débit passe de 37,6 à 46,1 pages par seconde
    (+23~\%). C'est le réglage le plus rentable.
  \item \textbf{Redis} pour les sessions et le cache applicatif : sans effet mesurable sur un
    serveur unique (le cache sur disque local est aussi rapide) ; il devient nécessaire avec
    plusieurs serveurs web.
  \item \textbf{Processeurs} : le débit croît avec le nombre de c\oe urs, la base n'étant pas le
    goulot (15~\% du temps au plus).
\end{itemize}

Les chiffres « après optimisation » du tableau de montée en charge combinent les optimisations du
module et les connexions persistantes : +27 à +35~\% de débit, temps de réponse réduits d'environ
20~\%.

\section{Optimisations supplémentaires (version 1.0.2)}

Une seconde série d'optimisations, chacune mesurée seule, dans l'ordre où elle a été ajoutée,
avec un protocole court (paliers de 50, 100 et 200 utilisateurs, 60~s chacun) et un profil CPU.
La mesure de référence a été refaite après un redémarrage de la machine.

\begin{center}
\begin{tabularx}{\linewidth}{@{}X r r r r@{}}
\toprule
\textbf{Étape (cumulée)} & \textbf{50} & \textbf{100} & \textbf{200} & \textbf{gain} \\
\midrule
Référence (version 1.0.1, connexions persistantes) & 44,4 & 47,2 & 46,3 & -- \\
Optimisations du module (ci-dessous) & 48,2 & 49,2 & 48,7 & +5 à +9~\% \\
APCu comme cache applicatif de Moodle & 53,5 & 55,5 & 57,7 & +11 à +18~\% \\
JIT de PHP 8.3 (\code{opcache.jit=tracing}) & 56,5 & 60,2 & 61,9 & +5 à +8~\% \\
OPcache sans vérification des fichiers & 60,7 & 59,2 & 62,0 & $\approx$ +2~\% \\
\bottomrule
\end{tabularx}
\end{center}
\noindent Pages servies par seconde, selon le nombre d'utilisateurs simultanés.

\subsection{Dans le module}
\begin{itemize}
  \item \textbf{Synthèse enseignant, liste « Validation »} : une seule requête couvre toutes les
    promotions, chacune avec ses propres droits ; la base trie et pagine, et seules les 40 lignes
    affichées sont mises en forme. 101 $\to$ 87~ms (143~ms avant la version 1.0.1).
  \item \textbf{Synthèse, promotions actives} : les six activités liées sont lues en une requête ;
    module de cours et contexte viennent des informations de cours déjà en cache de Moodle.
    7,6 $\to$ 3~ms, sans mise en cache des droits (le retrait d'un droit reste immédiat).
  \item \textbf{Pages secondaires} : la description de l'activité n'est plus mise en forme pour
    être aussitôt masquée ; 4 à 6~ms de moins par page.
  \item \textbf{Inscription directe depuis le catalogue} : l'année d'étude se choisit sur la ligne
    de l'EP, sans page de confirmation. Une ruée de 300 étudiants qui s'inscrivent au même instant
    (\code{ruee.py}) est servie en 11,6~s au lieu de 15,3~s (médiane par étudiant 9,2~s au lieu de
    13,0~s ; 900 pages au lieu de 1\,200 ; 300 inscriptions réussies dans les deux cas).
\end{itemize}

\subsection{Dans le déploiement}
\begin{itemize}
  \item \textbf{APCu} garde le cache applicatif de Moodle (chaînes de langue, configuration,
    informations de cours) en mémoire partagée au lieu de le relire sur disque à chaque page :
    c'est le gain le plus net de cette série. Le fichier de langue du module (près de
    1\,000 chaînes) en profite directement.
  \item \textbf{JIT} : le rendu des gabarits et le code de Moodle, très consommateurs de
    processeur, s'exécutent plus vite.
  \item \textbf{OPcache sans vérification des fichiers} : gain faible, dans le bruit de mesure.
  \item Avec APCu ou sans vérification des fichiers, il faut redémarrer PHP-FPM après chaque mise
    à jour du code ou mise à niveau faite en ligne de commande.
  \item \emph{Ajout ultérieur.} Le test des interactions à 1\,000 étudiants a montré qu'avec
    APCu, juste après un redémarrage de PHP-FPM ou une purge des caches, des connexions
    simultanées échouent sur le verrou du cache des activités des cours (\code{core/coursemodinfo}),
    qu'APCu ne sait pas tenir : ce cache doit rester sur le magasin fichier
    (rapport \emph{Test des interactions à 1\,000 étudiants}).
\end{itemize}

\subsection{Résultat}

\begin{center}
\begin{tabularx}{\linewidth}{@{}r R R R R R@{}}
\toprule
\textbf{Utilisateurs} & \textbf{Premier test} & \textbf{Version 1.0.1} & \textbf{Version 1.0.2} &
  \textbf{médiane 1.0.2} & \textbf{erreurs} \\
\midrule
25  & 37,3 & 47,4 & 58,9 & 0,42 s & 0 \\
50  & 36,5 & 45,6 & 59,8 & 0,83 s & 0 \\
100 & 37,2 & 46,8 & 60,6 & 1,67 s & 0 \\
200 & 37,9 & 49,7 & 59,8 & 3,48 s & 0 \\
300 & 42,1 & 51,2 & 61,1 & 5,24 s & 0 \\
400 & 44,1 & 53,1 & 60,9 & 7,33 s & 0 \\
\bottomrule
\end{tabularx}
\end{center}

Pages par seconde. Le débit d'un serveur de 4 c\oe urs passe d'environ 40 à 60 pages par seconde
(+40 à +60~\% depuis le premier test), et le temps de réponse médian diminue d'autant. Une
endurance de 10 minutes à 150 utilisateurs avec tous ces réglages tient 54 à 56 pages par seconde,
sans erreur et sans dérive de la mémoire.

\section{Endurance}

Trente minutes à 150 utilisateurs simultanés sans temps de réflexion, au-delà de la
saturation (\code{stress.py endurance}), soit environ 80\,000 pages servies :

\begin{center}
\begin{tabularx}{\linewidth}{@{}l R R R R R R@{}}
\toprule
\textbf{Fenêtre} & \textbf{pages/s} & \textbf{médiane} & \textbf{95\,\%} & \textbf{99\,\%} & \textbf{erreurs} &
  \textbf{mémoire} \\
\midrule
0--5 min   & 44,4 & 3,33 s & 3,67 s & 4,17 s & 0 & 2\,638 Mo \\
5--10 min  & 44,2 & 3,33 s & 3,85 s & 4,49 s & 0 & 2\,673 Mo \\
10--15 min & 43,5 & 3,41 s & 3,81 s & 4,21 s & 0 & 2\,748 Mo \\
15--20 min & 45,3 & 3,26 s & 3,70 s & 4,15 s & 0 & 2\,776 Mo \\
20--25 min & 46,5 & 3,18 s & 3,57 s & 3,96 s & 0 & 2\,780 Mo \\
25--30 min & 43,1 & 3,44 s & 3,80 s & 4,26 s & 0 & 2\,778 Mo \\
\bottomrule
\end{tabularx}
\end{center}

Aucune dérive : débit et temps de réponse constants, mémoire stabilisée après dix minutes (les
caches se remplissent puis cessent de croître), 17 connexions à la base (les 16 processus PHP et
la surveillance), aucune erreur ni verrou restant.

\section{Traitements lourds simultanés}

Pendant que 50 utilisateurs ouvrent leurs pages en boucle (\code{lourds.py}), sont lancés au même
instant :

\begin{center}
\begin{tabularx}{\linewidth}{@{}X r r@{}}
\toprule
\textbf{Traitement} & \textbf{Durée} & \textbf{Résultat} \\
\midrule
Restauration d'une archive de 60\,000 crédits dans une activité vide (ligne de commande) & 20 s &
  59\,550 crédits, 327 Mo \\
Export CSV détaillé de deux promotions (48\,000 lignes chacune) & 18 et 22 s & 7,8 Mo chacun \\
Livret PDF de deux promotions (200 EP, livret conservé effacé) & 41 et 48 s & 1,3 Mo chacun \\
Export complet de trois promotions par trois membres de la DEVE & 61 à 74 s & 3,9 à 8,3 Mo \\
Émission groupée de 300 suppléments au diplôme (tâche ad hoc) & 113 s & 300 émis, 157 Mo \\
\bottomrule
\end{tabularx}
\end{center}

\begin{itemize}
  \item Tous les traitements aboutissent, sans erreur, alors que le processeur est déjà saturé
    par la charge interactive : leurs durées sont celles d'un serveur occupé (l'export complet
    prend 10~s sur un serveur au repos).
  \item La mémoire reste bornée : 310~Mo au plus pour un processus PHP, 2,3~Go pour l'ensemble des
    processus PHP, 3,2~Go pour la machine.
  \item Les utilisateurs interactifs ne voient aucune erreur ; leurs temps de réponse s'allongent
    modérément (95\,\% des pages en moins de 2,1~s au lieu de 1,4~s, débit de 38,8 au lieu de
    45,6 pages par seconde).
  \item L'émission groupée tourne en tâche de fond à environ 4 suppléments par seconde : une
    promotion de 1\,600 étudiants demande moins de 7 minutes.
\end{itemize}

\section{Capacité}

Un serveur de 4 c\oe urs réglé comme recommandé sert environ 60 pages par seconde. Un utilisateur
réel ouvre une page toutes les 10 à 15 secondes : le serveur suffit pour 600 à 900 personnes
actives \emph{au même moment}, ce qui couvre largement l'usage courant d'une école (quelques dizaines de personnes
connectées en même temps).

Le seul moment exigeant est l'\textbf{ouverture des inscriptions} au catalogue, si toute une
promotion se connecte à la même minute. Le test de ruée (rapport \emph{Test de concurrence})
et les paliers ci-dessus montrent que le serveur ne commet alors aucune erreur et ne surréserve
aucune place : il fait patienter. Avec l'inscription directe de la version 1.0.2, 300 étudiants
qui s'inscrivent au même instant sont tous inscrits en 12~s ; pour 1\,000, comptez une
quarantaine de secondes au pire, et 8 c\oe urs divisent ce délai par deux.
Étaler l'ouverture des inscriptions par promotion (périodes d'inscription des EP) est l'autre
réponse, sans matériel supplémentaire.

\section{Conclusion}

Sous un stress bien au-delà de l'usage réel -- jusqu'à 400 utilisateurs sans temps de réflexion,
une demi-heure en continu, et les traitements les plus lourds lancés ensemble -- le module ne
commet aucune erreur, ne perd aucune requête, ne dérive pas en mémoire et garde des données
cohérentes. Au-delà de la saturation, le serveur fait patienter au lieu de céder.

La limite est celle du socle Moodle sur le matériel ; le module n'y prend que quelques pour cent.
Les optimisations du module (versions 1.0.1 et 1.0.2 : index, requêtes uniques paginées par la
base, calculs une seule fois, inscription directe) et du déploiement (connexions persistantes,
APCu, JIT) ont porté le débit d'un serveur de 4 c\oe urs d'environ 40 à 60 pages par seconde
(+40 à +60~\%), avec des temps de réponse réduits d'autant.

Les scripts sont dans \code{docs/audit/banc/} : \code{stress.py}, \code{profil.py},
\code{lourds.py}, \code{ruee.py}, \code{charge_setup.php} et \code{synthese_setup.php}.
\end{document}
